\begin{abstract}
法律判決檢索同時面臨長文本、稀疏引用訊號與高度領域化語言三個挑戰。THUIR@COLIEE 2023 的 Table~3 顯示：當 SAILER 與 BM25 都被限制為 512 tokens 時，SAILER 的 F1 高於 BM25；但當 BM25 不受最大長度限制時，其 F1 又反而高於受截斷的 SAILER。這促使本研究提出核心假設：長文本截斷，可能是 dense retrieval 在法律判決檢索中無法穩定超越 lexical retrieval 的主要原因之一。基於此假設，我們選用 ModernBERT-base 作為長文本 encoder，並在美國判決書語料上以 masked language modeling 進行法律領域繼續預訓練，再於 COLIEE 2025 與 2026 Task~1 上進行對比式 supervised fine-tuning。

受限於 NVIDIA RTX 4080 SUPER 16GB，在 \texttt{query:positive:negatives=1:1:15} 的設定下，8192 tokens 會直接造成 CUDA out-of-memory，因此本研究正式實驗採用 4096 tokens。美國判決書 continued pretraining 亦受 GPU 與時間成本限制，約訓練 60 小時後仍只覆蓋不到半個 epoch。實驗結果顯示四點結論。第一，query 表示方式比預期更重要：在 2026 valid 上，若使用整篇判決書作為 query（\texttt{processed}），SAILER 與 ModernBERT DAP+SFT 的 F1@5 分別為 0.1213 與 0.1650；若改用 THUIR 風格的 reference sentence extraction（\texttt{processed\_new}），則下降為 0.1022 與 0.1038。第二，ModernBERT 若未經 supervised fine-tuning 幾乎無法直接用於本任務，未 fine-tune 的原始模型在 2026 valid 上僅有 0.0142 的 F1@5，而探索階段的 DAP-only 設定在 2025 valid 也只有 0.0116。第三，依照類似 THUIR@COLIEE 2023 的靜態 BM25 hard negative 流程，模型在 COLIEE 2025 valid 上只有約 0.00515 的 F1，幾乎等於隨機猜測；改為每個 epoch 依模型相似度動態抽樣後，驗證表現提升超過一個數量級。第四，加入 scope filtering、LightGBM learning-to-rank 與 post-processing 後，FLNLP 最終提交在 COLIEE 2026 Task~1 取得 0.2935 的 F1，於 22 隊中排名第 7。整體而言，本研究說明：長文本處理、query 設計、負樣本選取、以及後段排序整合，是法律判決 dense retrieval 能否成功的四個關鍵因素。
\end{abstract}
